% texte.tex — f175r, vv. 1044–1081 (conversion initiale des travaux antérieurs)
\vers[n=1044, indent=2em, findent=1.2em]{}{}{ὃς ἐμὲ κροκόεντ᾽ ἐνέδυσεν·}{me revêtit de cette crocote,}
\vers[n=1045, indent=2em, findent=1.2em]{}{}{ἐπὶ δὲ τοῖσδε τόδ᾽ ἀνέπεμψεν}{puis de surcroît me dépêcha vers ce}
\vers[n=1046, indent=2em, findent=1.2em]{}{}{ἱερόν, ἔνθα γυναῖκες.}{temple, au milieu de femmes.}
\vers[n=1047, indent=2em, findent=1.2em]{}{}{ἰώ μοι μοίρας}{Hélas!}
\vers[n=1047, indent=2em, findent=1.2em]{}{}{ἄτεγκτε δαίμων·}{démon implacable de mon destin !}
\vers[n=1048, indent=2em, findent=1.2em]{}{}{ὦ κατάρατος ἐγώ·}{Ô moi le maudit !}
\vers[suite]{ΜΝ.}{Mnésiloque}{τίς ἐμὸν οὐκ ἐπόψεται}{Qui, en présence de mes misères, ne tournerait les yeux}
\vers[suite]{ΜΝ.}{Mnésiloque}{πάθος ἀμέγαρτον ἐπὶ κακῶν παρουσίᾳ;}{sur cette inenviable souffrance ?}
\vers[n=1050, indent=2em, findent=1.2em]{}{}{εἴθε με πυρφόρος αἰθέρος ἀστὴρ—}{Puisse l'Étoile dévorante de l'éther}
\vers[n=1051, indent=2em, findent=1.2em]{}{}{τὸν βάρβαρον ἐξολέσειεν.}{anéantir un infortuné !}
\vers[n=1052, indent=2em, findent=1.2em]{}{}{οὐ γὰρ ἔτ᾽ ἀθανάταν φλόγα λεύσσειν}{Car il ne m'est plus doux de voir la lumière immortelle,}
\vers[n=1053, indent=2em, findent=1.2em]{}{}{ἐστὶν ἐμοὶ φίλον, ὡς ἐκρεμάσθην,}{suspendu que je suis,}
\vers[n=1054, indent=2em, findent=1.2em]{}{}{λαιμότμητ᾽ ἄχη δία μένων αἰόλαν}{attendant comme terme d'un supplice jugulant}
\vers[n=1055, indent=2em, findent=1.2em]{}{}{νέκυσιν ἐπὶ πορείαν.}{un trajet rapide vers les morts.}
\vers[n=1056]{ΕΥ.}{Euripide}{χαῖρ᾽ ὦ φίλη παῖ· τὸν δὲ πατέρα Κηφέα}{(Dans la coulisse et sans être vu). Salut, chère enfant. Mais ton père Cepheus}
\vers[n=1057]{}{}{ὅς σ᾽ ἐξέθηκεν ἀπολέσειαν οἱ θεοί.}{qui t'exposa, que les dieux le confondent.}
\vers[n=1058]{ΜΝ.}{Mnésiloque}{σὺ δ᾽ εἶ τίς ἥτις τοὐμὸν ᾤκτιρας πάθος;}{Qui es-tu, toi qui compatis à mon tourment ?}
\vers[n=1059]{ΕΥ.}{Euripide}{Ἠχὼ, λόγων ἀντῳδὸς ἐπικοκκάστρια,}{Je suis Écho, la rediseuse goguenarde,}
\vers[n=1060]{}{}{ἥπερ πέρυσιν ἐν τῷδε ταὐτῷ χωρίῳ}{celle qui, l'an passé, en ce même lieu,}
\vers[n=1061]{}{}{Εὐριπίδῃ καὐτὴ ξυνηγωνιζόμην.}{prêtai mon aide à Euripide.}
\vers[n=1062]{}{}{ἀλλ᾽ ὦ τέκνον σὲ μὲν τὸ σαυτῆς χρὴ ποιεῖν,}{Songe, mon enfant, à tenir bien ta partie}
\vers[n=1063]{}{}{κλαίειν ἐλεινῶς.}{et à gémir lamentablement.}
\vers[suite]{ΜΝ.}{Mnésiloque}{σὲ δ᾽ ἐπικλαίειν ὕστερον.}{Et toi, à gémir de plus belle après moi.}
\vers[n=1064]{ΕΥ.}{Euripide}{ἐμοὶ μελήσει ταῦτά γ᾽. ἀλλ᾽ ἄρχου λόγων.}{Je n'y manquerai pas. Commence.}
\vers[n=1065, indent=3em, findent=1.7999999999999998em, metre={an.}]{ΜΝ.}{Mnésiloque}{ὦ νὺξ ἱερὰ}{Ô Nuit sainte,}
\vers[n=1066, indent=3em, findent=1.7999999999999998em]{}{}{ὡς μακρὸν ἵππευμα διώκεις}{quelle longue carrière tu fournis,}
\vers[n=1067, indent=3em, findent=1.7999999999999998em]{}{}{ἀστεροειδέα νῶτα διφρεύουσ᾽}{en guidant ton char sur le dos constellé}
\vers[n=1068, indent=3em, findent=1.7999999999999998em]{}{}{αἰθέρος ἱερᾶς}{de l'éther sacré,}
\vers[n=1068, indent=3em, findent=1.7999999999999998em]{}{}{τοῦ σεμνοτάτου δι᾽ Ὀλύμπου;}{à travers le très auguste Olympe.}
\vers[suite]{ΕΥ.}{Euripide}{†δι᾽ Ὀλύμπου;†}{Olympe.}
\vers[n=1070, indent=3em, findent=1.7999999999999998em]{ΜΝ.}{Mnésiloque}{τί ποτ᾽ Ἀνδρομέδα περίαλλα κακῶν}{Pourquoi seule Andromède reçut-elle tant de maux}
\vers[n=1071, indent=3em, findent=1.7999999999999998em]{}{}{μέρος ἐξέλαχον;}{en partage ?}
\vers[suite]{ΕΥ.}{Euripide}{μόρος ἔξελ᾽ ἄχος.}{En partage.}
\vers[n=1072, indent=3em, findent=1.7999999999999998em]{ΜΝ.}{Mnésiloque}{θανάτου τλήμων}{Malheureuse que la mort...}
\vers[suite]{ΕΥ.}{Euripide}{†θανάτου τλήμων†.}{Malheureuse que la mort.}
\vers[n=1073, indent=3em, findent=1.7999999999999998em]{ΜΝ.}{Mnésiloque}{ἀπολεῖς μ᾽ ὦ γραῦ στωμυλλομένη.}{Tu me tueras, vieille, avec ton babil.}
\vers[n=1074, indent=3em, findent=1.7999999999999998em]{ΕΥ.}{Euripide}{στωμυλιομενής.}{Ton babil.}
\vers[n=1075, indent=3em, findent=1.7999999999999998em]{ΜΝ.}{Mnésiloque}{νὴ Δί᾽, ὀχληρά γ᾽ εἰσήρρηκας}{Par Zeus, ton ingérence m'excède}
\vers[n=1076, indent=3em, findent=1.7999999999999998em]{}{}{λίαν.}{au dernier point.}
\vers[suite]{ΕΥ.}{Euripide}{†λίαν†.}{Au dernier point.}
\vers[n=1077, indent=3em, findent=1.7999999999999998em]{ΜΝ.}{Mnésiloque}{ὦγάθ᾽ ἔασόν με μονῳδῆσαι,}{Mon brave, laisse-moi chanter ma monodie.}
\vers[n=1078, indent=3em, findent=1.7999999999999998em]{}{}{καὶ χαριεῖ μοι παῦσαι.}{Tu m'obligeras : cesse !}
\vers[suite]{ΕΥ.}{Euripide}{᾽παΰσαι.}{(Criard :) C'est ce... que j'allais dire.}
\vers[n=1079, indent=3em, findent=1.7999999999999998em]{ΜΝ.}{Mnésiloque}{βάλλ᾽ ἐς κόρακας.}{Va-t'en aux corbeaux.}
\vers[suite]{ΕΥ.}{Euripide}{βάλλε σύ κόρακας.}{Va-t'en aux corbeaux.}
\vers[n=1080, indent=3em, findent=1.7999999999999998em]{ΜΝ.}{Mnésiloque}{†τί τὸ κακόν†;}{Qu'est-ce qui t'arrive ?}
\vers[suite]{ΕΥ.}{Euripide}{†τί τὸ κακόν;†}{Qu'est-ce qui t'arrive ?}
\vers[suite]{ΜΝ.}{Mnésiloque}{ληρεῖς.}{Tu veux rire.}
\vers[suite]{ΕΥ.}{Euripide}{†ληρεῖς†.}{Tu veux rire.}
\vers[n=1081, indent=3em, findent=1.7999999999999998em]{ΜΝ.}{Mnésiloque}{οἴμωζ᾽.}{Va te promener..}
\vers[suite]{ΕΥ.}{Euripide}{†οἴμωζ᾽†.}{Va te promener..}
\vers[suite]{ΜΝ.}{Mnésiloque}{ὀτότυζ᾽.}{Foin de toi.}
\vers[suite]{ΕΥ.}{Euripide}{†ὀτότυζ᾽†.}{Foin de toi.}
